Además de seleccionar un modelo por producto, se quiso evaluar si combinar dos modelos aporta algún valor en comparación con usar uno solo. Este análisis se agregó de manera independiente para el sistema Thary, ya que ninguna de las 4 implementaciones de referencia combina modelos. La idea viene de la literatura de combinación de pronósticos \cite{bates1969combination, granger1984improved}, que se explica en la Sección 2.1.

Un aspecto importante es que esta combinación no participa en la predicción que usa Thary en producción ni en la selección del modelo campeón, sino que solo se muestra en el panel técnico del administrador como un análisis complementario.

Se combinan las predicciones de XGBoost y Media Móvil porque, de los tres modelos candidatos, son los que obtuvieron mejores resultados en la validación (Tabla \ref{tab:resumen_cv}). Lo primero es calcular el peso $w$ que se le da a XGBoost, ya que a la Media Móvil le corresponde $1 - w$. Esto se hace con la fórmula de mínimos cuadrados explicada en la Sección 2.1, en la función \texttt{\_peso\_minimos\_cuadrados}:

\Needspace{10\baselineskip}
\captionof{listing}{Cálculo del peso de la combinación (función \texttt{\_peso\_minimos\_cuadrados}).}\label{lst:combinacion_peso}
\begin{minted}{python}
def _peso_minimos_cuadrados(y, pred_a, pred_b):
    diferencia = pred_a - pred_b
    denominador = float(np.sum(diferencia ** 2))
    if denominador == 0:
        return 0.5
    w = float(np.sum((y - pred_b) * diferencia) / denominador)
    return min(max(w, 0.0), 1.0)
\end{minted}

Aquí \texttt{y} es la demanda real, \texttt{pred\_a} es la predicción de XGBoost y \texttt{pred\_b} la de Media Móvil, y el cálculo del peso es la misma fórmula presentada en la Sección 2.1, aplicada sobre todas las filas del backtest. Hay dos casos especiales. El primero es que si las dos predicciones son idénticas, el denominador es cero, y como no hay nada que decidir entre un modelo y otro, el peso es 0.5, lo que además evita dividir entre cero. El segundo es que el resultado se recorta entre 0 y 1 con \texttt{min} y \texttt{max}, para que la combinación siempre sea un promedio ponderado, ya que si la demanda real queda fuera del rango entre las dos predicciones, la fórmula daría un peso mayor a 1 o menor a 0.

Para saber si la combinación realmente mejora los resultados, se compara contra XGBoost solo y Media Móvil sola sobre las mismas filas del backtest. Un aspecto importante es que el peso de cada mes se calcula sin usar ese mes, sino solo con los demás meses del backtest:

\Needspace{10\baselineskip}
\captionof{listing}{Cálculo del peso de cada mes sin usar ese mes (función \texttt{evaluar\_combinacion}).}\label{lst:combinacion_evaluacion}
\begin{minted}{python}
for mes in np.unique(meses):
    es_mes = meses == mes
    resto = ~es_mes
    if resto.any():
        w = _peso_minimos_cuadrados(y[resto], pred_xgb[resto], pred_mm[resto])
    else:
        w = 0.5

    pesos_por_mes[str(mes)] = round(w, 4)
    combinada[es_mes] = w * pred_xgb[es_mes] + (1 - w) * pred_mm[es_mes]
\end{minted}

Esto se hace porque si el peso se calculara con los mismos datos contra los que después se mide el error, el peso ya habría "visto" las respuestas de antemano y la combinación saldría artificialmente buena, que es el mismo principio de no hacer trampa que se explicó en la validación del modelo. Además, como los productos se juntan en un solo cálculo, hay miles de filas para calcular cada peso aunque cada producto tenga pocos meses. Si solo existe un mes, no quedan datos para calcular el peso y se usa 0.5.

Como resultado, la función devuelve las métricas (MAE, RMSE, Mediana del Error, R\textsuperscript{2}, MAPE y WAPE) de las tres opciones: Media Móvil sola, XGBoost solo y la combinación. También devuelve el peso final calculado con todos los meses, y cuánto mejora la combinación respecto al mejor de los dos modelos solos, en porcentaje de MAE y de RMSE. Un valor positivo significa que la combinación es mejor, y uno negativo que es peor. Esta función cuenta con 7 pruebas automatizadas, descritas en la Sección 3.4.1.